% ============================================================
\section{Results}
\label{sec:results}

Table~\ref{tab:main} reports the end-to-end matrix, each metric as a base/SFT
column pair. Appendix~\ref{app:full} gives all twelve model$\times$prompt cells.

\begin{table}[htbp]
\centering
{\tabstyle\setlength{\tabcolsep}{0.8pt}
\begin{tabular*}{\linewidth}[t]{@{\extracolsep{\fill}}l l c@{\hspace{2pt}}c c@{\hspace{2pt}}c c@{\hspace{2pt}}c c@{\hspace{2pt}}c@{}}\toprule
& & \multicolumn{2}{c}{\shortstack{\textbf{Refusal}\\\textbf{F1}}} & \multicolumn{2}{c}{\shortstack{\textbf{Behaviour}\\\textbf{Adherence}}} & \multicolumn{2}{c}{\shortstack{\textbf{Factual}\\\textbf{Grounding}}} & \multicolumn{2}{c}{\shortstack{\textbf{Single-Truth}\\\textbf{Recall}}}\\
\cmidrule(lr){3-4}\cmidrule(lr){5-6}\cmidrule(lr){7-8}\cmidrule(lr){9-10}
\textbf{Model} & \textbf{Prompt} & \textbf{base} & \textbf{\textsc{sft}} & \textbf{base} & \textbf{\textsc{sft}} & \textbf{base} & \textbf{\textsc{sft}} & \textbf{base} & \textbf{\textsc{sft}}\\\midrule
\multirow{2}{*}{Qwen2.5-7B} & strict & 0.429 & \textbf{0.911} & 0.811 & 0.780 & 0.617 & \textbf{0.838} & 0.561 & \textbf{0.687}\\
 & minimal & 0.031 & \textbf{0.881} & 0.737 & \textbf{0.768} & 0.031 & \textbf{0.837} & 0.757 & 0.686\\
\addlinespace[2pt]
\multirow{2}{*}{Qwen2.5-32B} & strict & 0.367 & \textbf{0.905} & 0.789 & \textbf{0.821} & 0.851 & 0.812 & 0.733 & \textbf{0.749}\\
 & minimal & 0.015 & \textbf{0.911} & 0.747 & \textbf{0.809} & 0.108 & \textbf{0.822} & 0.773 & 0.738\\
\addlinespace[2pt]
\multirow{2}{*}{Llama-3.1-8B} & strict & 0.255 & \textbf{0.926} & 0.721 & \textbf{0.805} & 0.705 & \textbf{0.852} & 0.634 & \textbf{0.727}\\
 & minimal & 0.015 & \textbf{0.934} & 0.743 & \textbf{0.795} & 0.074 & \textbf{0.857} & 0.710 & \textbf{0.714}\\
\addlinespace[2pt]
\multirow{2}{*}{Mistral-7B} & strict & 0.352 & \textbf{0.826} & 0.660 & \textbf{0.793} & 0.624 & \textbf{0.796} & 0.563 & \textbf{0.644}\\
 & minimal & 0.000 & \textbf{0.839} & 0.727 & \textbf{0.786} & 0.079 & \textbf{0.792} & 0.700 & 0.627\\
\bottomrule\end{tabular*}}

\caption{\textbf{End-to-end results on the 736-example held-out benchmark.} Within each pair the benchmark, prompts, decoding policy and evaluator are identical. Only the LoRA adapter differs. Bold marks a fine-tuned value that exceeds its baseline. Differences below the $\pm0.034$ noise floor (\S\ref{sec:setup}) are uninformative. We show the two endpoints of the instruction ladder. The intermediate runtime condition is in Appendix~\ref{app:full}.}
\label{tab:main}
\end{table}

\begin{table}[htbp]
\centering
{\tabstyle\setlength{\tabcolsep}{4pt}
\begin{tabular*}{\linewidth}[t]{@{\extracolsep{\fill}}l l rrrr@{}}\toprule
\textbf{Model} & \textbf{Method} & \textbf{Ref.\,F1} & \shortstack{\textbf{Behaviour}\\\textbf{Adherence}} & \shortstack{\textbf{Factual}\\\textbf{Grounding}} & \shortstack{\textbf{Single-Truth}\\\textbf{Recall}}\\\midrule
\multirow{3}{*}{Llama-3.1-8B} & CoN (prompt-only) & 0.000 & 0.632 & 0.609 & 0.658\\
 & CoT (few-shot) & 0.302 & 0.716 & 0.635 & 0.678\\
 & \textsc{trace-sft} (ours) & \textbf{0.926} & \textbf{0.805} & \textbf{0.852} & \textbf{0.727}\\
\addlinespace[2pt]
\multirow{3}{*}{Mistral-7B} & CoN (prompt-only) & 0.000 & 0.685 & 0.522 & 0.676\\
 & CoT (few-shot) & 0.525 & 0.657 & 0.417 & 0.528\\
 & \textsc{trace-sft} (ours) & \textbf{0.826} & \textbf{0.793} & \textbf{0.796} & \textbf{0.644}\\
\addlinespace[2pt]
\multirow{3}{*}{Qwen2.5-7B} & CoN (prompt-only) & 0.000 & 0.727 & 0.666 & 0.719\\
 & CoT (few-shot) & 0.402 & \textbf{0.781} & 0.734 & \textbf{0.724}\\
 & \textsc{trace-sft} (ours) & \textbf{0.911} & 0.780 & \textbf{0.838} & 0.687\\
\bottomrule\end{tabular*}}

\caption{\textbf{External prompting comparators} on baseline models; trace-tuned rows use the strict prompt, so every row receives a well-specified instruction. Both comparators reach usable answer content but neither acquires the answer/refuse decision.}
\label{tab:comparators}
\end{table}

\subsection{Refusal recall improves substantially}
\label{sec:refusal}

The largest and cleanest effect is the answer/refuse decision. Refusal F1
separates completely: every baseline configuration scores $\leq0.429$ and every
trace-tuned configuration $\geq0.826$, a mean improvement of $+0.709$. This is
not merely over-caution: decision accuracy also rises in every cell. The
benchmark is $82.6\%$ answerable, so answering everything scores $0.826$ for
free. No baseline configuration exceeds that floor by more than the $\pm0.034$
noise interval, while every trace-tuned cell reaches $0.927$--$0.976$.
Behaviour Adherence also improves, but unevenly, in an order that tracks
baseline instruction-following (Table~\ref{tab:main}; recipes differ by model
family, Appendix~\ref{app:limitations-expanded}).
Appendix~\ref{app:full} gives the per-cell recall and accuracy figures,
including the Llama-3.1-8B case, where decision accuracy masks near-zero
refusal recall.

\subsection{Learnt behaviour is prompt-invariant}
\label{sec:internalization}

Removing the prompt scaffolding is close to a null operation for fine-tuned
models and close to catastrophic for baseline ones. The minimal
condition is a sparse instruction, forty words that never name the
trace, the taxonomy or the citation format. A model that still produces them
is therefore reproducing a learned mapping rather than following the prompt.
Component by component, no trace-tuned model varies by more than the
$\pm0.034$ noise floor across the ladder. Baseline spreads are
$0.240$--$0.398$ on refusal F1 and $0.620$--$0.823$ on factual grounding.
Under the minimal prompt, baseline grounding collapses toward zero while
trace-tuned models hold steady (Table~\ref{tab:internalization}). The result also reproduces on
citation measurements built after and independently of CATS.
Figure~\ref{fig:internal} plots the full ladder. Appendix~\ref{app:figs} gives
the remaining figures.

\subsection{What trace supervision does \emph{not} buy}
\label{sec:whatdoesnot}

Two components show no consistent gain. Single-truth recall is positive in only
$7$ of $12$ cells (sign test $p=0.77$): fine-tuned models are not recovering
target answers more often. Factual grounding is not directly
comparable across run types, because it is averaged over
$\textit{fg\_n}=608+\textit{missed refusals}$. A baseline model that rarely
refuses is therefore additionally scored on ${\sim}130$ items whose evidence is
insufficient by construction. Bounding FG on the like-for-like population, the
intervals overlap in $6$ of $12$ cells, every one of them a strict or runtime
cell (Table~\ref{tab:fgdenom}; see Appendix~\ref{app:deconfound}). Trace supervision
therefore raises grounding when the instruction is sparse, but under a
well-specified prompt no consistent advantage is established. Given a
well-specified prompt, a baseline model already produces competitive grounded
answers. What it lacks is the judgment to decide whether to answer at all, and
the ability to retain any of this without being told the procedure. That is what
trace supervision supplies, which is why these negative results are predictions
of our account rather than exceptions to it.
Appendix~\ref{app:deconfound} gives the per-cell bounds and the
instruction-volume effect for baseline models.

\subsection{Mechanism-matched oracle ablations}
\label{sec:oracle}

Each oracle condition (\S\ref{sec:setup}) replaces one stage of the trace with ground truth
and leaves the rest of the pipeline unchanged. This isolates what that stage
contributes as a diagnostic intervention, never a deployment setting. Read as a single aggregate, their effect is small and the analysis stops there.
Read \emph{per metric}, with metrics the intervention should not affect acting
as a specificity control, the result is considerably stronger. Table~\ref{tab:oracle} and
Figure~\ref{fig:oraclefig} give the per-metric effect of each oracle, on which
the paragraphs below draw. Appendix~\ref{app:oracle} discusses them. \textbf{The conflict label moves behaviour and only behaviour.} Supplying the
gold Stage-2 label improves Behaviour Adherence ($+0.013$, $8/9$ cells,
$p=0.039$) (Table~\ref{tab:oracle}). \textbf{The document notes move grounding and only grounding.} Supplying gold
Stage-1 notes improves every grounding and citation metric, each positive in
\emph{every} cell (Table~\ref{tab:oracle}). Behaviour shows a
small spillover, expected since correct document verdicts also constrain the
conflict inference. The answer/refuse decision is untouched, since GR is
upstream of both interventions. \textbf{Interpretation.} That the effects add ($|\Delta_{\text{both}}-
(\Delta_{\text{conflict}}+\Delta_{\text{notes}})|\leq0.012$) is itself a
construct-validity argument for the decomposition. For \emph{baseline} models
every oracle instead \emph{hurts}, catastrophically so for grounding under the
conflict label ($-0.238$), a second internalization result.
Appendix~\ref{app:oracle} gives the per-metric residuals and the baseline
breakdown.

\subsection{External prompting comparators}
\label{sec:comparators}

We compare against two prompting-only strategies applied to baseline models,
neither fine-tuned and both on the same benchmark and evaluator. These are few-shot
Chain-of-Thought \citep{wei2022cot} and Chain-of-Note-style retrieval notes
\citep{yu2024chainofnote}. Both reach usable answer content (Table~\ref{tab:comparators}), but neither
acquires the answer/refuse decision. Chain-of-Note
\emph{never abstains} (Ref.\,F1$=0$ for every model, Table~\ref{tab:comparators}),
landing exactly on the answer-everything floor. Few-shot Chain-of-Thought
abstains inconsistently rather than rarely, from $32/128$ to $85/128$ across
models. Structured prompting can therefore elicit conflict-aware \emph{answers}
without supplying conflict-aware \emph{judgment}, the distinction
\S\ref{sec:intro} draws. Appendix~\ref{app:citation} gives the per-model
abstention spread and Chain-of-Note's citation hygiene, the poorest of any
method examined. \textbf{Citation behaviour.} Trace supervision changes citation \emph{sufficiency}, not citation precision: precision is near-ceiling for every model regardless of tuning, while the recall side moves sharply. Appendix~\ref{app:citation} gives the seven-metric analysis, the definitions and the per-cell values.

\subsection{Answer-only supervision comparison}
\label{sec:answeronly}

The comparators above are prompt-only. This one is a matched SFT comparator.
We fine-tune Llama-3.1-8B and Qwen2.5-7B on the identical data and QLoRA
configuration as the trace-supervised runs, but with the trace removed from
the supervision target. Appendix~\ref{app:answeronly} gives the exact target format. Because data, base model and optimiser are
held fixed, any remaining gap is attributable to the trace itself rather than
to additional supervision volume. Under the minimal prompt, trace supervision
exceeds this answer-only baseline on every metric for both models, with the
largest gap on Refusal Precision (Table~\ref{tab:aodelta}, Appendix~\ref{app:answeronly}).
Answer Quality (Appendix~\ref{app:aggregate}), the
geometric-mean combination of Factual Grounding and Single-Truth Recall used
elsewhere in the paper, favours the trace as well.
Appendix~\ref{app:answeronly} gives the complete per-metric deltas and the recipe
in full.

\section{Validating the Evaluator}
\label{sec:human}

A judged metric is only as trustworthy as its judge, so we validated the CATS
committee against an independent human study. Four reviewers double-judged 300
responses spanning both model families, both training conditions, and all five
conflict types. Grounded Refusal is deterministic and therefore outside this
exercise. Only the three judged components are at stake. \textbf{Factual Grounding
and Single-Truth Recall are reliable} ($\kappa=0.78$ and $0.79$ against humans).
\textbf{Behaviour Adherence is not} ($\kappa=0.41$), with the weakness
concentrated in Complementary information and Misinformation, where human
reviewers also disagree most. We therefore report conflict-conditioned behaviour
qualitatively and do not interpret individual BA margins. Appendices
\ref{app:human} and~\ref{app:allagreement} give the study design and coverage
caveats. Any reading of a BA number should be checked against them.

\section{Conclusion}
\label{sec:conclusion}

Trace supervision transforms the answer/refuse decision
and makes conflict-aware behaviour survive removal of the instruction describing
it. It does \emph{not} improve answer writing, because that deficit was never
there to fix. The three stages behave as separable constructs, so the trace
is a measurement instrument rather than a formatting convention. Three questions remain open. The residual error after both oracles is small but unattributed, and separating
Stage-3 synthesis failure from annotation noise and judge strictness needs a
targeted audit. Behaviour Adherence is the weakest instrument we report, so
conflict-behaviour claims should stay qualitative until it improves. Misinformation remains the smallest class in the benchmark. Expanding it without
perturbation-based synthesis is an open data-collection problem. The practical consequence is a separation of concerns. A retrieval system that
answers well but cannot decide \emph{when} to answer is not a weaker version of
one that can. The two differ in a capability that no prompt in our ladder
supplied, and that supervising the reasoning did. Building such a system means
supervising the judgment, and evaluating it means scoring that judgment apart
from the answer it produces.

% ============================================================
\section*{Limitations}
\label{sec:limitations}

\paragraph{Per-document verdicts rest on a small, non-blind check.} The
Stage-1 verdicts that gate Factual Grounding eligibility come from the
annotation committee. A targeted review of 100 query--document pairs (three
reviewers, 300 ratings) found $92\%$ agreement ($\kappa=0.88$) with the
human-majority verdict, against human--human agreement of $\kappa=0.83$.
Reviewers saw the committee verdict before rating, so anchoring may inflate
this figure, and no independent gold-standard label exists.

\paragraph{CATS is evidence-bounded, not truth-bounded.} It scores whether the
retrieved set justifies an answer, not whether the answer is true. The
grounding criterion is permissive by design, so our grounding numbers are an
upper bound relative to strict entailment.

\paragraph{Behaviour Adherence is our least reliable instrument}
($\kappa=0.41$ against humans, $\alpha=0.46$ agreement among judges). A priority-weighted vote rule makes disputed
judgments closer to a single-judge call than the panel size suggests. We
therefore report conflict-behaviour findings qualitatively rather than by
margin. Appendix~\ref{app:limitations-expanded} gives six further limitations:
model families, aggregate properties, recipe asymmetry, annotation quality,
review coverage, and comparator scope.

% ============================================================
\section*{Reproducibility Statement}

We release the training/validation split, the held-out benchmark, the
stagewise annotation prompts, the end-to-end inference prompts, the CATS
judge prompts, and the per-run evaluation outputs from which every number in
this paper is computed (Appendix~\ref{app:annprompts}). Appendix~\ref{app:sft}
gives the QLoRA hyperparameters, the two training recipes, and the
checkpoint-selection criterion for every fine-tuned model.
Appendix~\ref{app:inference} gives the deterministic decoding protocol
(temperature $0$, fixed continuation budgets, sentinel-based stopping) used
for every generation reported in the paper. Appendix~\ref{app:annotation}
gives the formal annotation aggregation rule and Appendix~\ref{app:catsprompts}
gives the complete CATS judge prompts, so both the released labels and the
evaluation metric can be independently reproduced. Appendix~\ref{app:data}
documents the retrieval protocol, the source-balanced sampling rule, and the
per-source composition of the benchmark.

% ============================================================
\section*{Ethics Statement}

\paragraph{Data provenance and licensing.} We derive the released resource from
publicly available research datasets and from live web retrieval. We preserve
source provenance in every record. The dataset cards identify the
upstream collection for each source family, so downstream users can consult
the relevant licences before redistributing source-derived material. We store
retrieved snippets with their source URL and, where available, a timestamp.
Because live retrieval is time-varying, readers should treat the release as a
snapshot, not a reproducible query against a current index.

\paragraph{Human annotators.} Members of the research group, not crowdworkers,
carried out the human review. We retain reviewer identifiers in
the review artifacts for auditability, but they carry no demographic information.
We make no claims about recruitment, compensation, or demographics that our
records do not support.

\paragraph{Intended use and risks.} We intend the resource for research on
evidence-grounded generation, abstention calibration, and conflict-aware
evaluation. First, conflict labels and answerability judgments
describe the \emph{supplied} evidence, not world truth. Readers must not treat
them as factual adjudications of the underlying questions, several of which concern
contested scientific or medical topics. Second, the source-quality field is a
coarse URL-whitelist heuristic retained as annotation metadata. It is not a
credibility ranking of publishers, and researchers should not use it as one.

% ============================================================
\section*{AI Use Statement}

\paragraph{Required disclosure.} We used frontier language models for
\textit{synthetic data generation}: they produced the synthetic annotations in
our dataset and served as the \textit{LLM-as-a-judge} evaluators of model
performance; the methodology for both is specified in \S\ref{sec:data} and
\S\ref{sec:cats}. No other required-disclosure task (developing the
theoretical framework, formulating claims, designing the methodology or
experiments, or implementing the fine-tuning and evaluation code) was
performed by generative AI; these are the authors' own.

\paragraph{Recommended disclosure.} We also used generative AI tools to edit
this manuscript for readability and length. This included tightening prose,
checking citation formatting and cross-reference integrity, and compressing
the main text to the submission page limit.

\paragraph{Accountability.} All experimental design, analysis, and claims are
the authors' own. We reviewed every edited passage against the source
results, and we take responsibility for the final content of this work,
including text, claims, and artifacts produced with the aid of generative AI.
